\section{项目背景与目标}
\label{sec:intro}
\subsection{项目背景}
车载摄像头为驾驶员状态分析提供了连续的面部画面。表情能够反映面部动作与交互状态，而同一画面中的五官比例、局部纹理和脸部外观又包含身份信息。因此，本项目计划研究一条先匿名、后分析的处理路线，在削弱身份关联的同时保留表情判断所需的信息。

课程项目拟以图片、视频和摄像头作为输入，通过 Python 桌面应用展示处理前后的对照。系统不以遮挡整张脸为终点，而是将匿名结果继续送入表情模型，最终输出表情类别、置信度与模型分析速度，完成匿名化和识别的联合演示。

\subsection{研究问题}
首先需要确定可以修改哪些面部信息。眼睛、眉毛、鼻部和嘴部既与身份有关，也承担表情表达，编辑时应改变外观而保留动作关系。其次需要明确处理边界：原始图像可用于人脸定位和对照显示，表情分类器只接收匿名灰度脸。最后需要考虑连续视频，使结果能跟随当前人脸，避免逐帧识别的孤立波动影响阅读。

\subsection{项目目标}
项目拟实现统一人脸匿名化模块、匿名图表情识别模块和桌面应用。匿名化限定于面部内部，保持头发、背景和座舱环境；表情识别覆盖七类基本表情，训练和运行使用一致的匿名输入定义；桌面应用提供样本对照、视频分析和整体指标三个功能区。

评价将围绕三个目标展开：降低自建身份模型对匿名脸的原身份命中，保持匿名图的表情可识别性，获得可观察的连续视频分析速度。项目将把准确性放在优先位置，同时设计适合持续视频演示的计算配置。

\subsection{项目范围}
本项目计划使用通用表情图像与车载近红外样本，完成数据整理、模型训练、联合评价和桌面集成。身份模型只作为项目内部的离线评价工具，不承担实际驾驶员身份业务。预期交付包括可运行源代码、运行资源说明、项目计划书、最终报告和演示材料。
